%\evensidemargin = 0.46cm
%\oddsidemargin = 0.46cm
%\topmargin = -1.54cm
%\headheight = 1cm
%\headsep = 0.5cm
%\textwidth = 14cm
%\textheight = 25.7cm
%\footskip = 1.25cm
%\setlength{\parindent}{0ex}
%\setlength{\parskip}{2ex}
\pagestyle{headings}
%\marginparsep = 0.5cm
%\marginparwidth = 2cm
%\hbadness=10000

\def\beforepreface{
\pagenumbering{roman}
\pagestyle{plain}
\thispagestyle{empty}%
\pagebreak[4]%\newpage
\setcounter{page}{1}
}

\def\prefacesection#1{%
\parskip 0pt
\chapter*{#1}
\addcontentsline{toc}{chapter}{#1}}
\parskip 0pt

\def\afterpreface{%\newpage
\parskip 0pt
\addcontentsline{toc}{chapter}{Table of Contents}
\addtolength{\parskip}{-0.75ex}
\tableofcontents
\addtolength{\parskip}{+0.75ex}
\cleardoublepage
\pagenumbering{arabic}
\pagestyle{fancy}
\renewcommand{\footrulewidth}{0.4pt}}



\makeatletter
\def\thickhrulefill{\leavevmode \leaders \hrule height 1pt\hfill \kern \z@}
\renewcommand{\maketitle}{\begin{titlepage}%
    \let\footnotesize\small
    \let\footnoterule\relax
    \parindent \z@
    \reset@font
%    \null
    \vskip 10\p@
    \hbox{\mbox{%
        \hspace{4pt}%
        \includegraphics[width=4cm]{epsimages/logohighschool.eps}%
        \hspace{4pt}
        }%
      \vrule depth 0.9\textheight%
      \mbox{\hspace{2em}}
      \vtop{% %%%%%%%%%%%%%%%%%%
        \vskip 40\p@
        \begin{flushleft}
          \Large \@author \par
        \end{flushleft}
        \vskip 80\p@
        \begin{flushleft}
          \huge \bfseries \@title \par
        \end{flushleft}
        \vskip 80\p@
        \begin{center}
          \large \bfseries Version 0.5.1 \par
          \large \bfseries \@date \par
        \end{center}
        \vfil
        }}
    \null
  \end{titlepage}%
  \setcounter{footnote}{0}%
}\makeatother


%===========================================================================%
% Here are the locally defined macros which are designed to make the FHSST
% authors' lives just that little bit more bearable ;-)

% Note To Self (comes up parenthesised and red)
\newcommand{\nts}[1]{(\textcolor{red}{NOTE TO SELF: #1})}
\newcommand{\query}[1]{(\textcolor{red}{Query: #1})}
\newcommand{\ohm}{\ensuremath{\Omega}}
\newcommand{\eohm}{\,\Omega}
\newcommand{\eN}{\,\rm{N}}                %m in text
\newcommand{\emm}{\,\rm{m}}                %m in text
\newcommand{\ep}{\,\ekg \cdot \mbox{\ms}}                %m/s in text
\newcommand{\es}{\,\mathrm{s}}                %s in equation
\newcommand{\ekg}{\,\mathrm{kg}}                %kg in equation
\newcommand{\eJ}{\,\mathrm{J}}                %J in equation
\newcommand{\eA}{\,\mathrm{A}}                %A in equation
\newcommand{\eV}{\,\mathrm{V}}                %J in equation
\newcommand{\eW}{\,\mathrm{W}}                %W in equation
\newcommand{\ms}{m$\cdot$s$^{-1}$}                %m/s in text
\newcommand{\mss}{m$\cdot$s$^{-2}$}                %m/s in text
\newcommand{\ems}{\,\mathrm{m \cdot s^{-1}}}            %m/s in equation
\newcommand{\emss}{\,\mathrm{m \cdot s^{-2}}}            %m/s in equation
\newcommand{\qstart}[1]{#1\begin{enumerate}}
\newcommand{\qend}{
\end{enumerate}}
\newcommand{\px}{$x$}                % position x, in text
\newcommand{\py}{$y$}                % position y, in text
\newcommand{\edx}{\Delta x}        % displacement dx, in text
\newcommand{\dx}{$\edx$}            % displacement dx, in text
\newcommand{\edy}{\Delta y}            % displacement dy, in text
\newcommand{\dy}{$\edy$}            % displacement dy, in text
\newcommand{\edt}{\Delta t}            % delta time dt, in text
\newcommand{\dt}{$\edt$}            % delta time dt, in text
\newcommand{\vel}{$v$}                % velocity
\newcommand{\kph}{km$\cdot$hr$^{-1}$}    %km/h in text
\newcommand{\momen}{\vec{p}}            %momentum
\newcommand{\kener}{KE}                            %kinetic energy
\newcommand{\poten}{PE}                            %kinetic energy
\newcommand{\degree}{^{\circ}}
\newcommand{\ie}{{\em i.e.~}}
\newcommand{\eg}{{\em e.g.~}}
\newcommand{\cf}{{\em c.f.~}}
\newcommand{\resp}{{\em resp.~}}
\newcommand{\etc}{{\em etc.~}}
\newcommand{\nb}{{\em n.b.~}}
\def\deg{$^\circ$}
\newcommand{\ud}{\mathrm{d}}
\newcommand{\comment}[1]{}



%===========================================================================%
% New style colours

\definecolor{siyavula-gray}{RGB}{242,242,242}
\definecolor{siyavula-dark-gray}{RGB}{205,204,204}

%===========================================================================%
% Classico font as per Peter JOhn
\renewcommand*\sfdefault{uop}
\renewcommand*\familydefault{\sfdefault}

%===========================================================================%


\newcommand{\practiceinfo}[0]{
\par\vspace*{2em}
\raisebox{-0.7em}{\includegraphics[scale=0.5]{../../icons/aplus.ps}} More practice 
\phantom{xx}\raisebox{-0.7em}{\includegraphics[scale=0.5]{../../icons/video.ps}} video solutions 
\phantom{xx}\raisebox{-0.7em}{\includegraphics[scale=0.5]{../../icons/help.ps}} or help at www.everythingmaths.co.za
\vspace*{2em}
}

% This is an environment for conveniently writing out the syllabus for a
% particular section. using \item for each relevant entry from the syllabus
\newenvironment{syllabus}{\color{blue}The syllabus requires:\begin{itemize}}{\end{itemize}}

\newcounter{workedexamplecounter} % physics worked example counter
\setcounter{workedexamplecounter}{1}
\newlength{\saveparindentlength}
\setlength{\saveparindentlength}{\parindent}
% Essay counter :)
\newcounter{essaycounter} % physics worked example counter
\setcounter{essaycounter}{1}
% My attempt to define an essay environment
% Will take 3 arguments
% 1 - Essay title
% 2 - Authors name and details block
% 3 - About the author block

\newenvironment{essay}[4]%
{\begin{center}\begin{pspicture}(-5,0)(5,0.01)
 \psline[linecolor=red]{-}(5,0)(7,0)\psline[linecolor=red]{-}(7,0)(7,-1)
 \psline[linecolor=red]{-}(-7,0)(-5,0)\psline[linecolor=red]{-}(-7,0)(-7,-1)
\end{pspicture}\end{center}
\vspace{.5cm}\setlength{\parindent}{5pt}
 \textbf{\textit{ Essay \arabic{essaycounter}} \\ #1 \small{by #2}}%
 \addtocounter{essaycounter}{1}%
 \par\vspace{.5cm} #4 \\\par\vspace{0.25cm}\textit{\textbf{Author: #2}}\\\vspace{0.5cm}%
\small{#3}\\ %
}%
{%
\begin{center}\begin{pspicture}(0,0)(5,0.01)%
 \psline[linecolor=red]{-}(0,0)(5,0)%
\end{pspicture}\end{center}%
\setlength{\parindent}{\saveparindentlength}}

\def\essay[#1]{\pagebreak \addcontentsline{toc}{chapter}{\textit{Essay \arabic{essaycounter}}: #1}%
 \textbf{\Large\textit{Essay \arabic{essaycounter}}: #1\addtocounter{essaycounter}{1}}%
}
\def\essayauthor[#1]{\vspace{.5cm}\textit{\textbf{Author: #1}}}
\def\essayauthorblurb[#1]{\par\vspace{.5cm}{\sffamily #1}}
\def\essaytitle[#1]{\vspace{0.3cm}{\textbf{\Large #1}}\par\vspace{0.3cm}}

% My attempt to define a worked example environment
\newcounter{stepcounter}
\setcounter{stepcounter}{1}

% can we deprecate the use of this, and use \westep instead
\newcommand{\westep}[1]
{\par \vspace{3mm}%\pagebreak[1]
\leftskip20pt
\noindent
\textit{Step \arabic{stepcounter}} : \textit{\textbf{#1}}
\par\noindent
\leftskip50pt
\addtocounter{stepcounter}{1}}

\ifx\setlinejoinmode\undefined
 \newcommand{\setlinejoinmode}[1]{}
\fi
\ifx\setlinecaps\undefined
 \newcommand{\setlinecaps}[1]{}
\fi
% This way define your own fonts mapping (for example with ifthen)
\ifx\setfont\undefined
 \newcommand{\setfont}[2]{}
\fi

\newcommand{\pscompass}%
{\scalebox{1} % Change this value to rescale the drawing.
{
\begin{pspicture}(-1.5,-1.561875)(1.5,1.561875)
\pspolygon[fillstyle=solid,fillcolor=black](0,0.75)(0,0)(-0.15,0.15)
\pspolygon(0,0.75)(0,0)(0.15,0.15)
\pspolygon(0.75,0)(0,0)(0.15,-0.15)
\pspolygon[fillstyle=solid,fillcolor=black](0.75,0)(0,0)(0.15,0.15)
\pspolygon[fillstyle=solid,fillcolor=black](0,-0.75)(0,0)(0.15,-0.15)
\pspolygon(0,-0.75)(0,0)(-0.15,-0.15)
\pspolygon(-0.75,0)(0,0)(-0.15,0.15)
\pspolygon[fillstyle=solid,fillcolor=black](-0.75,0)(0,0)(-0.15,-0.15)
\uput[u](0,0.75){N}
\uput[d](0,-0.75){S}
\uput[l](-0.75,0){W}
\uput[r](0.75,0){E}
\end{pspicture} 
}



}

\newcommand{\psrings}%
{%
\pscircle*[linewidth=0.5pt](0cm,0cm){.01}
\pscircle[linewidth=.5pt,linecolor=gray](0cm,0cm){.25}%
\pscircle[linewidth=.5pt,linecolor=gray](0cm,0cm){.5}%
\pscircle[linewidth=.5pt,linecolor=gray](0cm,0cm){0.75}%
\pscircle[linewidth=.5pt,linecolor=gray](0cm,0cm){1}%
}
\newcommand{\psintfact}%
{%
\rput(0,0){\begin{tabular}{c}\textbf{Interesting}\\ \textbf{Fact}%
\end{tabular}}
}
\newcommand{\psfact}%
{%
\begin{pspicture}(-.5,-.5)(.5,.5)
 \psintfact
 \PstLens[LensSize=0.5,LensRotation=-60,LensMagnification=1.5](-.04,-.1){\psintfact}
\end{pspicture}
}
\newcommand{\psflag}%
{%
\scalebox{0.3} % Change this value to rescale the drawing.
{\begin{pspicture}(0,-1.92)(3.0744452,1.92)
\definecolor{color792b}{rgb}{0.6,0.6,0.6}
\psline[linewidth=0.04cm](0.14559634,1.06)(0.9655963,-1.9)
\psline[linewidth=0.04cm](0.24559633,1.06)(1.0855963,-1.88)
\rput{13.182965}(0.26058337,-0.015092337){\psellipse[linewidth=0.04,dimen=outer](0.19559634,1.12)(0.19,0.1)}
\psline[linewidth=0.04cm](0.94559634,-1.88)(1.0655963,-1.86)
\psbezier[linewidth=0.04,fillstyle=solid,fillcolor=color792b](0.86559635,1.42)(1.2655964,1.64)(0.94559634,1.24)(1.1455964,1.2)(1.3455963,1.16)(2.4855964,1.9)(2.5255964,1.86)(2.5655963,1.82)(3.054445,0.5009174)(2.9255962,0.54)(2.7967474,0.5790826)(1.4855963,-0.3)(1.3855963,-0.16)(1.2855964,-0.02)(1.3455963,-0.08)(1.3455963,0.12)(1.3455963,0.32)(0.7855963,-0.28)(0.7055963,-0.34)(0.62559634,-0.4)(0.26559633,0.84)(0.28559634,0.86)(0.30559632,0.88)(0.26559633,0.94)(0.26559633,0.94)(0.26559633,0.94)(0.46559635,1.2)(0.86559635,1.42)
\end{pspicture}}
}
\newcommand{\psqmark}%
{%
\scalebox{0.2} % Change this value to rescale the drawing.
{\begin{pspicture}(0,-3.51)(3.98,3.53)
\psbezier[linewidth=0.04,fillstyle=solid,fillcolor=black](0.92,2.95)(1.22,3.37)(1.84,3.51)(2.44,3.43)(3.04,3.35)(3.96,2.93)(3.88,1.69)(3.8,0.45)(2.06,-0.27)(1.88,-1.09)(1.7,-1.91)(2.3,-2.09)(2.36,-2.15)(2.42,-2.21)(1.44,-1.99)(1.44,-1.15)(1.44,-0.31)(2.94,0.77)(2.84,1.69)(2.74,2.61)(1.36,3.07)(0.68,1.95)(0.0,0.83)(0.62,2.53)(0.92,2.95)
\pscircle[linewidth=0.04,dimen=outer,fillstyle=solid,fillcolor=black](2.09,-2.96){0.55}
\end{pspicture}}
}
\newcommand{\pstip}%
{%
\scalebox{0.2} % Change this value to rescale the drawing.
{\begin{pspicture}(0,-3.1533968)(2.7831953,3.1533968)
\definecolor{color2871}{rgb}{0.4,0.4,0.4}
\psbezier[linewidth=0.04,linecolor=color2871,fillstyle=solid,fillcolor=color2871](0.30975282,1.480703)(0.59950,0.346823)(1.4712996,-1.9338118)(1.9878281,-1.8548366)(2.5043566,-1.7758614)(2.18284,0.61323)(1.8401965,1.87331)(1.497551,3.13339)(0.019999,2.6145823)(0.30975,1.480703)
\rput{14.334663}(-0.565312,-0.6214416){\pscircle[linewidth=0.04,linecolor=color2871,dimen=outer,fillstyle=solid,fillcolor=color2871](2.18828,-2.55848){0.54}}
\end{pspicture}}
}
\newcommand{\psbulb}%
{%
\scalebox{0.2} % Change this value to rescale the drawing.
{\begin{pspicture}(0,-3.3013647)(4.908481,3.4294605)
\definecolor{color946b}{rgb}{0.6,0.6,0.6}
\psbezier[linewidth=0.04](3.2474866,-1.5642114)(2.4098418,0.36813208)(1.6433032,0.23539852)(1.6970811,0.6480823)(1.7508591,1.0607661)(2.1953533,0.77506554)(2.2262053,0.587854)(2.2570574,0.4006425)(1.5684397,0.80635864)(2.196053,1.1597726)(2.8236666,1.5131865)(2.8261008,0.56492805)(2.579102,0.8190769)(2.3321033,1.0732257)(2.7962835,1.2613873)(2.784706,1.2439616)(2.7731285,1.2265358)(2.950391,1.4443324)(3.1070092,1.2329777)(3.2636273,1.0216233)(2.798124,0.61800635)(3.7428837,-1.3030947)
\psbezier[linewidth=0.08](1.5593854,2.8649697)(2.7758908,3.3884606)(4.093773,2.5972388)(4.3687177,1.4005746)(4.643662,0.20391025)(3.7490253,-0.52576435)(4.290271,-1.442636)(4.831516,-2.3595076)(4.802724,-2.4393263)(4.7687063,-2.3669193)(4.7346883,-2.294512)(4.867481,-2.7182646)(4.2435164,-2.9893143)(3.6195517,-3.2603643)(3.4195778,-3.0228555)(3.4642856,-3.0239484)(3.5089936,-3.025041)(3.3216572,-3.002568)(2.9688413,-2.0634615)(2.6160254,-1.1243552)(1.6755583,-1.3425524)(0.83777916,-0.30250326)(0.0,0.73754585)(0.34288007,2.3414786)(1.5593854,2.8649697)
\psbezier[linewidth=0.06,fillstyle=solid,fillcolor=color946b](3.5100658,-3.0310266)(3.7100658,-3.2394476)(4.8300657,-2.849974)(4.730066,-2.3194475)(4.630066,-1.7889212)(4.550066,-1.9194475)(4.3100657,-1.4594476)(4.070066,-0.9994476)(2.810066,-1.5794476)(2.9700658,-2.0794475)(3.130066,-2.5794475)(3.310066,-2.8226056)(3.5100658,-3.0310266)
\psbezier[linewidth=0.05,fillstyle=solid,fillcolor=color946b](3.079205,-2.3223355)(3.1871026,-2.2979457)(3.9018052,-1.5315504)(3.9259355,-1.4160248)(3.9500659,-1.3004992)(4.010066,-1.2194476)(3.7197256,-1.3377068)(3.4293854,-1.4559659)(3.0753205,-1.9640672)(2.962319,-2.1313665)(2.8493176,-2.2986655)(2.9713078,-2.346725)(3.079205,-2.3223355)
\psbezier[linewidth=0.05,fillstyle=solid,fillcolor=color946b](3.2135053,-2.5354118)(3.3468883,-2.504099)(4.230406,-1.520147)(4.260236,-1.3718271)(4.290066,-1.2235073)(4.1889153,-1.1194476)(4.0053186,-1.271277)(3.8217223,-1.4231063)(3.208703,-2.0754423)(3.0690103,-2.2902327)(2.9293175,-2.505023)(3.0801222,-2.566725)(3.2135053,-2.5354118)
\psbezier[linewidth=0.05,fillstyle=solid,fillcolor=color946b](3.3335052,-2.735412)(3.4668882,-2.704099)(4.3504057,-1.720147)(4.3802357,-1.5718272)(4.410066,-1.4235073)(4.308915,-1.3194476)(4.1253185,-1.471277)(3.9417222,-1.6231064)(3.328703,-2.2754421)(3.1890101,-2.4902325)(3.0493176,-2.705023)(3.200122,-2.766725)(3.3335052,-2.735412)
\psbezier[linewidth=0.05,fillstyle=solid,fillcolor=color946b](3.4935052,-2.975412)(3.6268883,-2.944099)(4.5104055,-1.960147)(4.540236,-1.8118272)(4.570066,-1.6635073)(4.468915,-1.5594476)(4.285319,-1.711277)(4.1017222,-1.8631064)(3.488703,-2.5154421)(3.3490102,-2.7302325)(3.2093174,-2.945023)(3.3601222,-3.006725)(3.4935052,-2.975412)
\psbezier[linewidth=0.05,fillstyle=solid,fillcolor=color946b](3.6135051,-3.155412)(3.7468882,-3.1240988)(4.630406,-2.140147)(4.660236,-1.9918271)(4.690066,-1.8435073)(4.5889153,-1.7394476)(4.4053187,-1.891277)(4.221722,-2.0431063)(3.548703,-2.7754421)(3.46901,-2.9102325)(3.3893175,-3.045023)(3.480122,-3.1867251)(3.6135051,-3.155412)
\rput{31.819052}(-0.8005619,-2.576656){\psellipse[linewidth=0.07,dimen=outer,fillstyle=solid,fillcolor=color946b](4.119578,-2.692639)(0.76116425,0.3465934)}
\rput{25.244768}(-0.8473483,-2.0893452){\psellipse[linewidth=0.04,dimen=outer,fillstyle=solid,fillcolor=black](4.241361,-2.9366095)(0.27659246,0.1380608)}
\psbezier[linewidth=0.04](2.403768,2.3592985)(2.611458,2.299041)(2.721967,2.2683775)(3.0574007,1.9846212)(3.3928344,1.7008649)(3.2636173,1.9681277)(3.5846286,1.9902756)(3.90564,2.0124235)(3.2477834,2.7296348)(2.714062,2.59565)(2.1803405,2.4616654)(2.1960783,2.4195557)(2.403768,2.3592985)
\end{pspicture}}
}
\newcommand{\pspencil}%
{%
\scalebox{0.2} % Change this value to rescale the drawing.
{\begin{pspicture}(0,-2.6732397)(5.2826033,2.5843153)
\definecolor{color3983b}{rgb}{0.6,0.6,0.6}
\rput{44.4269}(1.5588851,-0.010889986){\psframe[linewidth=0.04,dimen=outer,fillstyle=solid,fillcolor=black](1.2927761,2.2532396)(0.29277614,1.5532397)}
\psbezier[linewidth=0.04](1.0711963,0.767678)(0.9788245,0.8230626)(2.0244849,1.4042634)(1.9913505,1.4501121)(1.9582161,1.495961)(4.1760063,-0.58293444)(4.514382,-1.123584)(4.8527575,-1.6642336)(5.150486,-2.583232)(5.2065444,-2.6125135)(5.2626033,-2.6417952)(4.336381,-2.2293487)(3.6961355,-1.8021653)(3.0558898,-1.3749818)(1.1635681,0.7122935)(1.0711963,0.767678)
\psbezier[linewidth=0.04,fillstyle=solid,fillcolor=black](4.7548966,-2.3972247)(4.9297223,-2.5004773)(5.151456,-2.6131008)(5.206459,-2.6331704)(5.261462,-2.6532397)(5.10783,-2.3825366)(5.0126433,-2.1965938)(4.9174566,-2.010651)(4.580071,-2.293972)(4.7548966,-2.3972247)
\psline[linewidth=0.04cm](1.3127761,0.99323976)(3.7127762,-1.3467603)
\psline[linewidth=0.04cm](1.6727761,1.2132398)(4.0727763,-1.1267602)
\psbezier[linewidth=0.04](4.0327764,-1.0867603)(4.292776,-1.3867602)(4.6527762,-1.1667602)(4.5327764,-1.2067603)
\psbezier[linewidth=0.04](3.612776,-1.2467602)(4.352776,-1.7467602)(4.1527762,-1.4667603)(4.0871353,-1.1067603)
\psbezier[linewidth=0.04](3.6927762,-1.3067603)(3.8927763,-1.7067603)(3.7527761,-1.9667603)(3.792776,-1.8067603)
\rput{41.58525}(1.5649645,0.18445837){\psellipse[linewidth=0.04,dimen=outer,fillstyle=solid,fillcolor=black](0.539593,2.1529272)(0.4898439,0.2609555)}
\psbezier[linewidth=0.04,fillstyle=solid,fillcolor=color3983b](1.3428409,2.242443)(1.6086246,2.2932398)(1.5615585,2.142183)(1.4889802,2.0629127)(1.4164017,1.9836423)(1.8161182,2.040834)(1.794309,1.7979672)(1.7724998,1.5551004)(2.352776,1.7629943)(1.9469734,1.4005489)(1.5411707,1.0381036)(2.0222244,1.4581892)(1.844831,1.3085054)(1.6674373,1.1588216)(1.937567,1.409145)(1.7413607,1.2381673)(1.5451542,1.0671896)(1.4235525,0.9589731)(1.1182237,0.7161064)(0.8128949,0.47323975)(0.90062475,1.0635847)(0.85651326,0.9810519)(0.81240183,0.8985192)(0.5511844,1.0472883)(0.616612,1.2239187)(0.6820397,1.4005489)(0.4606455,1.093056)(0.3767108,1.3343126)(0.29277614,1.5755692)(1.0770572,2.1916463)(1.3428409,2.242443)
\psbezier[linewidth=0.04](1.0027244,1.0147622)(1.1309127,1.3605639)(1.5727762,1.6332397)(1.7327762,1.6532397)
\psbezier[linewidth=0.04](0.88272446,1.1747621)(1.0109127,1.520564)(1.3527762,1.7132398)(1.5127761,1.7732397)
\psbezier[linewidth=0.04](0.7427244,1.3147621)(0.87091273,1.660564)(1.2127762,1.8532398)(1.3727762,1.9132397)
\rput{-141.02766}(10.480751,-1.0945002){\pstriangle[linewidth=0.04,dimen=outer,fillstyle=solid](5.046733,-2.6609855)(0.21336946,0.5188883)}
\end{pspicture}}
}

\newenvironment{TempEnv}[3]%
{% This is a picture at the beginning, a line or bar or something else
\begin{center}
   \begin{pspicture}(0,-0.2)(5,0.2)
      \psline[doubleline=true,linewidth=1pt]{-}(0,0)(5,0)
      \psline[doubleline=true,linewidth=1pt,border=3pt](2,0.2)(2.5,-0.2)(3,0.2)
   \end{pspicture}
\end{center}% End of opening picture
% this command puts a bar in the margin and sets the bar width
\setlength{\changebarwidth}{8pt}\cbstart
% Now we set margins
\begin{quotation}
\noindent
\vspace{.5cm}
\setlength{\parindent}{0pt}
%% Title and counter
\textbf{\textit{Worked Example \arabic{workedexamplecounter}: }
#1}\par
\textbf{Question:} #2\par
\textbf{Answer}\par #3%
\addtocounter{workedexamplecounter}{1}\setcounter{stepcounter}{1}%
\par%\vspace{.25cm}
}%
{\cbend
\end{quotation}\begin{center}\begin{pspicture}(0,-0.2)(5,0.2)
\psline[doubleline=true,linewidth=1pt]{-}(0,0)(5,0)
\psline[doubleline=true,linewidth=1pt,border=3pt](2,-0.2)(2.5,0.2)(3,-0.2)
\end{pspicture}\end{center}
\setlength{\parindent}{\saveparindentlength}}




\newcommand{\IFact}[1]{
\marginpar{%odd
\begin{minipage}[m]{0.225\paperwidth}
\begin{mdframed}[linecolor=black, linewidth=1.0pt,leftmargin=15pt, rightmargin=15pt]
\textbf{\textit{FACT}}
\end{mdframed}
\begin{mdframed}[backgroundcolor=lightgray, linecolor=lightgray,skipabove=-4pt, leftmargin=15pt, rightmargin=15pt]
\small #1
\end{mdframed}
\end{minipage}
}
}




%Defining our own aside environment!
\newcommand{\Extension}[2]
{%\vspace{3mm}
%\begin{quotation}
%\rput(-2.5,0){\psbulb}\noindent%
%\textit{Extension: #1}\\
%#2
%\end{quotation}
%\vspace{3mm}
\setlength{\saveparindentlength}{\parindent}
\setlength{\parindent}{0pt}

\ifdimless{\pagegoal-\pagetotal}{4em}%4em play with this 
{\clearpage\vspace*{-\topskip}}{} 

\begin{quotation}
\vspace*{25.6pt}
{  \fontsize{12}\bfseries\itshape
\begin{tikzpicture}[remember picture,overlay]
    \node[anchor=west ,text centered, xshift=-40pt]
      {\begin{tikzpicture}[remember picture, overlay]
        \draw[color=white, draw=black, line width=1pt] (0,0) rectangle
          (0.2\textwidth, 28pt);
       \node[anchor=west, yshift=14pt, xshift=4pt, text centered]
              { \bfseries ~ \color{black} Extension:};
       \end{tikzpicture}
    };
    
% now the white bar to the end of the paragraph
    \node[anchor=west, xshift=45pt]{
      \begin{tikzpicture}[remember picture, overlay]
        \draw[fill=white, draw=black, line width=1pt] (0,0) rectangle
          (0.8\textwidth-5pt,28pt);
         \node[anchor=west, yshift=14pt, xshift=0pt, text centered]
              { \textit{#1} };

       \end{tikzpicture}


      };
   \end{tikzpicture}

}

\begin{mdframed}[leftline=true, topline=false, bottomline=true, rightline=false, linecolor=siyavula-gray, linewidth=10pt]
\vspace{5pt}

#2

\end{mdframed}
\end{quotation}
\setlength{\parindent}{\saveparindentlength}








}







\newcommand{\Definition}[2]{
\begin{mdframed}[nobreak=true, linewidth=1,linecolor=black, skipabove=36pt]%
\vspace*{10pt}
\leftskip20pt
\rightskip20pt
{\Large\textsl{DEFINITION: } \textit{#1}}

\par\noindent
#2
\vspace*{1em}
\end{mdframed}
}



\newcommand{\HistoricalNote}[2]{
\vspace{3mm}
\begin{center}
\psshadowbox[framearc=0.2]{\begin{minipage}[t]{0.95\textwidth}
\textbf{Historical Note: #1}\par
#2
\end{minipage}}
\end{center}
\vspace{3mm}}


\mathchardef\ocomma="013B
\mathchardef\pcomma="613B
\mathcode`\,="8000
{\catcode`\,=\active
 \gdef,{\obeyspaces
 \futurelet\next\smartcomma}}
\def\smartcomma{\if\space\next
 \pcomma\else\ocomma\fi}

%the macro divrule#1#2 produces a rule the same width as
%#2 digits, indented by the width of #1 digits.
%~ is used as an active character to produce a space equal to a comma
% this is in order to align numbers properly
\newdimen\digitwidth
\settowidth\digitwidth{,}
%{\catcode`\~=\active
% \gdef~{\hspace{\digitwidth}}}
\def\divrule#1#2{%
 \noalign{\moveright#1\digitwidth%
 \vtop{\hrule width#2\digitwidth}}}

% Adds a reminder to put an example in.
\newcommand{\Ex}{(\textcolor{red}{ADD AN EXAMPLE HERE})}
\newcommand{\new}[1]{{\textit #1}}
\newcommand{\advanced}[1]{\nts{Advanced: #1}}
\def\deg{$^\circ$}

\newcommand{\equ}[2]{\begin{equation}#1\label{#2}\end{equation}}
\newcommand{\nequ}[1]{\begin{equation*}#1\end{equation*}}

\newcommand{\Tip}[1]{
\marginpar{%odd
\begin{minipage}[m]{0.225\paperwidth}
\begin{mdframed}[linecolor=black, linewidth=1.0pt, leftmargin=15pt, rightmargin=15pt]
\textbf{\textit{Tip}}
\end{mdframed}
\begin{mdframed}[backgroundcolor=lightgray, linecolor=lightgray,skipabove=-4pt,leftmargin=15pt, rightmargin=15pt]
\small #1
\end{mdframed}
\end{minipage}
}
}

\newcommand{\GlossEntry}[2]
{%\vspace{3mm}
%\begin{center}
\marginpar{
\textit{\textbf{#1} : #2}
}
%\newglossaryentry{#1}{name=#1,description=#2}
%\end{center}
%\vspace{3mm}
}

\newcommand{\MarginTip}[1]{
\marginpar{\begin{center}
\psshadowbox{\begin{minipage}[t]{\marginparwidth}%
 #1%
 \end{minipage}}
\end{center}}
}

\newcommand{\Activity}[3]{
%\vspace{3mm}
%\begin{center}
%\begin{pspicture}(0,-0.1)(10,0.1)\psline[linewidth=1pt]{-}(0,0)(10,0)
%\end{pspicture}
%  \begin{quotation} % or save, then modify margin
%     \textbf{Activity :: #1 : #2}\par #3%
%  \end{quotation} % or set margins back to normal
%\begin{pspicture}(0,-0.1)(10,0.1)\psline[linewidth=1pt]{-}(0,0)(10,0)
%\end{pspicture}
%\end{center}
%\vspace{3mmi}
\setlength{\saveparindentlength}{\parindent}
\setlength{\parindent}{0pt}

\ifdimless{\pagegoal-\pagetotal}{4em}%4em play with this 
{\clearpage\vspace*{-\topskip}}{} 

\begin{quotation}
\vspace*{25.6pt}
{  \fontsize{12}\bfseries\itshape
\begin{tikzpicture}[remember picture,overlay]
    \node[anchor=west ,text centered, xshift=-40pt]
      {\begin{tikzpicture}[remember picture, overlay]
        \draw[color=white, draw=black, line width=1pt] (0,0) rectangle
          (0.2\textwidth, 28pt);
       \node[anchor=west, yshift=14pt, xshift=4pt, text centered]
              { \bfseries ~ \color{black} Activity:};
       \end{tikzpicture}
    };
    
% now the white bar to the end of the paragraph
    \node[anchor=west, xshift=45pt]{
      \begin{tikzpicture}[remember picture, overlay]
        \draw[fill=white, draw=black, line width=1pt] (0,0) rectangle
          (0.8\textwidth-5pt,28pt);
         \node[anchor=west, yshift=14pt, xshift=0pt, text centered]
              { \textit{#2} };

       \end{tikzpicture}


      };
   \end{tikzpicture}

}

\begin{mdframed}[leftline=true, topline=false, bottomline=false, rightline=false, linecolor=siyavula-dark-gray, linewidth=10pt]
\vspace{5pt}

#3

\end{mdframed}
\end{quotation}
\setlength{\parindent}{\saveparindentlength}


}
% End of Activity


\newcommand{\Aim}[1]{\par\noindent\textbf{Aim:} \par\noindent #1}
\newcommand{\Apparatus}[1]{\par\noindent\textbf{Apparatus:} \par\noindent #1}
\newcommand{\Method}[1]{\par\noindent\textbf{Method:} \par\noindent #1}
\newcommand{\Results}[1]{\par\noindent\textbf{Results:} \par\noindent #1}
\newcommand{\Conclusions}[1]{\par\noindent\textbf{Conclusions:} \par\noindent
#1}



%EZ new counters for exercises, 
\newcounter{exercisecounter}
\setcounter{exercisecounter}{1}
% make it reset for new chapters
\makeatletter
\@addtoreset{exercisecounter}{chapter}
\makeatother

\newenvironment{eocexercises}[1]
{
\vspace*{1cm}
\eocexercisesheading
\addcontentsline{toc}{section}{End of chapter exercises}
%\hrule width \textwidth height 1pt\\
\begin{quotation}
\noindent
%\vspace{.5cm}
\setlength{\parindent}{0pt}
%\textbf{\textit{End of chapter exercises:}
%#1}
\par
}%
{%
\end{quotation}
\setlength{\parindent}{\saveparindentlength}
}


\newcommand{\Exercise}[2]{
\vspace*{0.5cm}
\exercisesheading
#2
}


\newtheorem{Theorem}{Theorem}
\newtheorem{Lemma}{This is a leemmma}

\newcommand{\mytheorem}[3]{
{\begin{Theorem}{#2}\label{#1}\end{Theorem}}
\textbf{Proof}:{#3}}

\newcommand{\schooltheorem}[6]{
{\begin{Theorem}{#2}:#3\label{#1}\end{Theorem}}
\textbf{Given}:{#4} \par  
\textbf{R.T.P.}:{#5} \par
\textbf{Proof}:{#6}
}



\newcounter{stepcounter}
\setcounter{stepcounter}{1}


\newcounter{workedexamplecounter}
\setcounter{workedexamplecounter}{1}

\newenvironment{wex}[3]
{\newlength{\myfloatsep}
\setlength{\myfloatsep}{\floatsep}
\setlength{\floatsep}{2mm}
\begin{quotation}
\noindent
\newlength{\oldparindent}
\setlength{\oldparindent}{\parindent}
\setlength{\parindent}{20pt}

\ifdimless{\pagegoal-\pagetotal}{4cm}%4em play with this 
{\clearpage\vspace*{-\topskip}}{} 


\begin{mdframed}[linewidth=1pt,linecolor=black, leftmargin=10pt, innertopmargin=8pt, innerbottommargin=6pt, skipbelow=0pt]
{\Large\textbf{\textit{Worked Example \arabic{workedexamplecounter}:}} } {\Large\textit{#1}}
\vspace*{0pt}
\end{mdframed}
\vspace*{-1em}
\begin{mdframed}[linecolor=siyavula-gray,backgroundcolor=siyavula-gray, leftmargin=10pt, skipbelow=8pt]
{\large\textbf{\textit{QUESTION}}} \vspace*{4pt} 
\hrule
\vspace*{18pt}
\noindent
\textit{\noindent #2}
\vspace*{28pt}
\par\noindent
\textbf{\textit{SOLUTION}}\vspace*{4pt}
\hrule
\vspace*{18pt}
\noindent
 #3%
\vspace*{28pt}
\hrule
\addtocounter{workedexamplecounter}{1}\setcounter{stepcounter}{1}%
%\vspace{.25cm}
}%
{%
\end{mdframed}
\end{quotation}
%\nopagebreak[4]
%\begin{pspicture}(0,0)(10,0.01)
% \psline[linecolor=gray,linewidth=3pt]{-}(0,0)(10,0)
% \psline[linecolor=gray]{-}(10,0)(10,2)
%\end{pspicture}
\setlength{\parindent}{\oldparindent}
\setlength{\floatsep}{\myfloatsep}
}


\newcommand{\MarginCompass}{\marginpar{\begin{center}\pscompass\end{center}}}

\newcommand{\todo}[1]{(\textcolor{red}{To Do: #1})}

%Chemistry Definitions
\def\beaker{\psline[linearc=7pt](0,2)(0,0)(1.5,0)(1.5,2)\psellipse(0.75,2)(0.8,0.1)}
\def\filledbeaker{\psline[linearc=7pt](0,2)(0,0)(1.5,0)(1.5,2)
\psline[fillstyle=solid,fillcolor=lightgray,linearc=7pt](0,1)(0,0)(1.5,0)(1.5,1)\psellipse(0.75,2)(0.8,0.1)}

\def\bunsen{
\pnode(0,0){A}\pnode(0,0.2){C}\pnode(1.4,0.2){D}
\pnode(0.6,0.4){E}\pnode(0.8,0.4){F}\pnode(0.7,0.4){EF}
\pnode(0.6,1.4){K}\pnode(0.8,1.4){L}
\pnode(0,0.6){G}\pnode(0.6,0.6){H}
\pnode(0,0.7){I}\pnode(0.6,0.7){J}
\pnode(0.5,0.75){J1}\pnode(0.9,0.95){J2}
\pnode(0.5,1.6){M}\pnode(0.9,2){N}
\pnode(0.55,1.4){O}\pnode(0.85,1.6){P}
\psframe(A)(D)
\psline(C)(E)\psline(D)(F)
\psellipse(EF)(0.1,0.05)
\psframe[fillstyle=solid,linestyle=none,fillcolor=white](E)(0.9,0.5)
\psline(E)(K)\psline(F)(L)
\psline(G)(H)\psline(I)(J)
\psframe(M)(N)\psframe(O)(P)
\psframe[fillstyle=solid,fillcolor=white](J1)(J2)
\pscircle(0.7,0.85){0.05}\pscircle(0.6,0.85){0.05}\pscircle(0.8,0.85){0.05}}

\def\stand{\psline(0,0)(0.5,2.5)\psline(0,2.5)(2.5,2.5)\psline(2.5,0)(2,2.5)}


%Physics Definitions
\def\rope{\psline(0,0.2)(5,.2)\psline(0,0)(5,0)\multirput(0,0)(0.2,0){25}{\psline(0,0)(0.2,0.2)}}

\def\cart{\psframe(0,0.4)(1,1)\pscircle(0.2,0.2){0.2}\pscircle(0.8,0.2){0.2}}

% Arrow for objects and images
\newpsobject{oi}{psline}{arrowsize=6pt, arrowlength=1.5, arrowinset=0, linewidth=2pt}
\psset{lensHeight=3,lensColor=lightgray}
\newpsobject{PrincipalAxis}{psline}{linewidth=0.5pt,linecolor=gray}

%Maths Definitions
\def\arrow{\psline(0,0)(0.2;-45)\psline(0,0)(0.2;45)}

\newcommand{\fline}[1]{\mathrm{#1}}

\psset{gridcolor=lightgray,RightAngleSize=0.2}

\def\andgate{\rput(1.2,0){
\psline(0.75,1)(0,1)(0,0)(0.75,0)
\psarcn(0.75,0.5){0.5}{90}{-90}
\psline(1.25,0.5)(2,0.5)
\psline(-0.75,0.3)(0,0.3)
\psline(-0.75,0.7)(0,0.7)
\uput[l](-0.75,0.7){A}
\uput[l](-0.75,0.3){B}
\uput[r](2,0.5){AB}}}

\def\nandgate{\rput(1.2,0){
\psline(0.75,1)(0,1)(0,0)(0.75,0)
\psarcn(0.75,0.5){0.5}{90}{-90}
\psline(1.25,0.5)(2,0.5)
\psline(-0.75,0.3)(0,0.3)
\psline(-0.75,0.7)(0,0.7)
\psdot[dotstyle=o](1.32,0.5)
\uput[l](-0.75,0.7){A}
\uput[l](-0.75,0.3){B}
\uput[r](2,0.5){$\overline{\rm{AB}}$}}}

\def\notgate{\rput(1.2,0){
\pspolygon(0.75,0.5)(0,1)(0,0)
\psdot[dotstyle=o](0.8,0.5)
\psline(0.8,0.5)(1.5,0.5)
\psline(-0.75,0.5)(0,0.5)
\uput[l](-0.75,0.5){A}
\uput[r](1.5,0.5){$\bar{\rm{A}}$}}}

\def\orgate{\rput(1.2,0.5){
\pscurve(-0.1,0.5)(0,0)(-0.1,-0.5)
\psline(-0.1,-0.5)(0.1,-0.5)
\pscurve(0.1,-0.5)(0.5,-0.4)(0.95,0)
\psline(-0.1,0.5)(0.1,0.5)
\pscurve(0.1,0.5)(0.5,0.4)(0.95,0)
\psline(-0.75,-0.3)(-0.05,-0.3)
\psline(-0.75,0.3)(-0.05,0.3)
\psline(0.95,0)(1.7,0)
\uput[r](1.7,0){A+B}
\uput[l](-0.75,0.3){A}
\uput[l](-0.75,-0.3){B}}}

\def\norgate{\rput(1.2,0.5){
\pscurve(-0.1,0.5)(0,0)(-0.1,-0.5)
\psline(-0.1,-0.5)(0.1,-0.5)
\pscurve(0.1,-0.5)(0.5,-0.4)(0.95,0)
\psline(-0.1,0.5)(0.1,0.5)
\pscurve(0.1,0.5)(0.5,0.4)(0.95,0)
\psline(-0.75,-0.3)(-0.05,-0.3)
\psline(-0.75,0.3)(-0.05,0.3)
\psline(0.95,0)(1.7,0)
\psdot[dotstyle=o](0.97,0)
\uput[r](1.7,0){$\overline{\rm{A+B}}$}
\uput[l](-0.75,0.3){A}
\uput[l](-0.75,-0.3){B}}}


% just remove sections form header - too long
\renewcommand{\sectionmark}[1]{\markright{}{}}







% Fancy headers 

\newcounter{shortcodecounter}
\setcounter{shortcodecounter}{1}




\usepackage[explicit]{titlesec}

\newcommand*\chapterlabel{}

\titleformat{\chapter}
  {\gdef\chapterlabel{}
   \normalfont\sffamily\fontsize{38}\bfseries\slshape}
  {\gdef\chapterlabel{\thechapter\ }}{0pt}
  {\begin{tikzpicture}[remember picture,overlay]
    \node[yshift=-4cm] at (current page.north west)
      {\begin{tikzpicture}[remember picture, overlay]
        \shade[left color=white, right color=gray] (0,0) rectangle
          (1.2\paperwidth,4cm);
        \node[anchor=west,text centered, yshift=2.5cm, xshift=1cm]
              {
              \begin{minipage}{1.2\textwidth}
              \raggedright
              \phantom{xxx}\color{black} #1
              \end{minipage}
              };
        \node[anchor=west,text centered, yshift=2.0cm, xshift=0.85\paperwidth]
              {\begin{minipage}{\textwidth}
              \color{white}\chapterlabel
              \end{minipage}
              };
      \end{tikzpicture}
      };
   \end{tikzpicture}
  }
\titlespacing*{\chapter}{1.5cm}{1cm}{1cm}



\newcommand*\sectionlabel{}
\titleformat{\section}
  {\gdef\sectionlabel{}
   \normalfont\sffamily\fontsize{22}\normalfont\itshape}
  {\gdef\sectionlabel{\thesection}}{0pt}
{
\begin{tikzpicture}[remember picture,overlay]
   \node[anchor=west, yshift=-1.5cm, xshift=-4.0cm]
      {\begin{tikzpicture}[remember picture, overlay]
        \draw[fill=siyavula-gray, color=siyavula-dark-gray, xshift=-\marginparwidth] (0,0) rectangle
          (1.5\paperwidth,2.0cm);
        \node[anchor=west,text centered ,yshift=1.0cm, xshift=\hoffset+0.45\marginparwidth]
              {
              \begin{minipage}{0.85\textwidth}
              \raggedright
              \color{black} #1
              \end{minipage}
              };
        \node[anchor=east,text centered, yshift=1.0cm, xshift=\hoffset+0.45\marginparwidth+1.1\textwidth]
              {\Large\color{black} \raisebox{-0.25em}{\includegraphics[scale=0.3]{../../icons/www.ps}}~P\arabic{shortcodecounter}};
       \end{tikzpicture}
      };
   \end{tikzpicture}
\addtocounter{shortcodecounter}{1}
}
\titlespacing*{\section}{1cm}{1cm}{2cm}




\newcommand*\subsectionlabel{}
\titleformat{\subsection}
  {\gdef\subsectionlabel{}
   \normalfont\sffamily\fontsize{18}\bfseries\slshape}
  {\gdef\subsectionlabel{\thesubsection}}{0pt}
  {\begin{tikzpicture}[remember picture,overlay]
   \node[anchor=west, yshift=-1.5cm, xshift=-0.5cm]
      {\begin{tikzpicture}[remember picture, overlay]
        \draw[fill=siyavula-gray, color=siyavula-gray] (0,0) rectangle
          (0.915\paperwidth,1.5cm);
        \node[anchor=west,text centered ,yshift=0.75cm, xshift=\hoffset+0.02\marginparwidth]
              { \begin{minipage}{0.85\textwidth}
                \color{black} #1
                \end{minipage}
                };
        \node[anchor=east,text centered, yshift=0.75cm, xshift=\hoffset+0.05\marginparwidth+1.1\textwidth]
              {\Large\color{black} \raisebox{-0.25em}{\includegraphics[scale=0.3]{../../icons/www.ps}}~P\arabic{shortcodecounter}};
       \end{tikzpicture}
      };
   \end{tikzpicture}
\addtocounter{shortcodecounter}{1}
  }
\titlespacing*{\subsection}{0cm}{1cm}{2cm}







\newcommand{\exercisesheading}{
\addtocounter{exercisecounter}{1}
 \par  \vspace*{2\baselineskip}
{  \fontsize{12}\bfseries\itshape
\begin{tikzpicture}[remember picture,overlay]
    \node[anchor=west ,text centered]
      {\begin{tikzpicture}[remember picture, overlay]
        \draw[color=white, draw=black, line width=1pt] (0,0) rectangle
          (0.3\textwidth, 28pt);
       \node[anchor=west, yshift=14pt, xshift=4pt, text centered]
              { \bfseries ~ \color{black} Exercise \chapterlabel~-~\arabic{exercisecounter}   };
       \end{tikzpicture}
    };
    
% now the gray bar to the end of the paragraph
    \node[anchor=west, xshift=0.335\textwidth]{
      \begin{tikzpicture}[remember picture, overlay]
        \draw[fill=white, color=white] (0,0) rectangle
          (0.655\textwidth,28pt);
       \end{tikzpicture}


      };
   \end{tikzpicture}

}
\par
%\vspace{\baselineskip}
}


\newcommand{\eocexercisesheading}{
\par \vspace*{3\baselineskip} 
{\normalfont \sffamily\fontsize{16}\bfseries
\begin{tikzpicture}[remember picture,overlay]
    \node[anchor=west, yshift=6\baselineskip, xshift=0., text centered]
      {\begin{tikzpicture}[remember picture, overlay]
        \draw[color=white, draw=black, very thick] (0,0) rectangle
          (0.3\textwidth,1.5cm);
       \node[anchor=west, yshift=0.75cm, text centered]
              { ~~  \color{black} Chapter \chapterlabel };
       \end{tikzpicture}
    };
    
% now the gray bar to the end of the paragraph
    \node[anchor=west, xshift=0.31\textwidth]{
      \begin{tikzpicture}[remember picture, overlay]
        \draw[fill=siyavula-gray, color=siyavula-gray] (0,0) rectangle
          (0.69\textwidth,1.5cm);
            \node[anchor=west, yshift=0.75cm, text centered]
              {   \color{black} End of Chapter Exercises };

       \end{tikzpicture}


      };
   \end{tikzpicture}

}
\par
\vspace{\baselineskip}
}



\newcommand{\summary}[1]{
\vspace{6\baselineskip} \par
{\normalfont \sffamily\fontsize{22}\bfseries
\begin{tikzpicture}[remember picture,overlay]
    \node[anchor=west, xshift=-0.4\paperwidth]
      {\begin{tikzpicture}[remember picture, overlay]
        \draw[fill=siyavula-gray, color=siyavula-gray] (0,0) rectangle
          (1.2\paperwidth,2cm);
        \node[anchor=west,text centered ,yshift=1.cm, xshift=0.3\paperwidth]
              {  \phantom{xxx} \color{black} Chapter \chapterlabel | Summary  };
       \end{tikzpicture}
      };
   \end{tikzpicture}

}
\par
\vspace{\baselineskip}
\addcontentsline{toc}{section}{Summary}

See the summary presentation (\raisebox{-0.2em}{\includegraphics[height=1em]{../../icons/video.ps}} Presentation: #1 at www.everythingmaths.co.za)

}
